\chapter{Uvod}
\label{ch:uvod}

\begin{izvor}
\textbf{Izvor:} \texttt{01\_Uvod.pdf}, stranice 1--79 (obrađeno u cijelosti).
\end{izvor}

\section{Što je optimizacija?}
\str{2--7}

\textbf{Etimologija.} Riječ \emph{optimizacija} potječe od latinske riječi
\textbf{„optimus”}, što znači \emph{najbolji}.

\begin{definicija}[opisna, sa slajda]
Optimizacija sadrži skup aktivnosti čiji je cilj pronaći \textbf{„najbolje
rješenje”}, tzv.\ \textbf{optimum}, uz poštivanje zadanih \textbf{ograničenja}.
\end{definicija}

Uočite da su u toj rečenici već sadržana \textbf{tri} ključna sastojka svakog
optimizacijskog problema, i vrijedi ih odmah imenovati jer ćemo se na njih
vraćati kroz cijeli kolegij:

\begin{enumerate}
  \item \textbf{Nešto se bira} --- postoje veličine koje mi možemo mijenjati
        (npr.\ dimenzije konstrukcije, količina proizvedene robe, iznosi struja).
  \item \textbf{Postoji mjerilo „najboljeg”} --- moramo imati broj koji svakom
        izboru pridružuje ocjenu (cijena, masa, potrošnja\dots), inače riječ
        „najbolje” nema značenja.
  \item \textbf{Ne smije se birati što god} --- izbori moraju zadovoljiti
        ograničenja (npr.\ ne možemo proizvesti više nego što tvornica može,
        dimenzija ne može biti negativna).
\end{enumerate}

\textbf{Optimizacija kao matematička disciplina.} Sa slajda:

\begin{quote}
Optimizacija = matematička teorija i algoritmi („matematička mašinerija”) ne
samo za donošenje najboljih odluka, već često i ključna za donošenje odluka
općenito (izabrati najbolje rješenje od svih mogućih \emph{vs.}\ da li uopće
postoji rješenje).
\end{quote}

Ta je zagrada važnija nego što se čini. Kod složenih inženjerskih problema
često uopće \textbf{nije očito postoji li ijedan izbor} koji zadovoljava sva
ograničenja. Pitanje „postoji li dopušteno rješenje?” je samo po sebi
optimizacijski problem --- i to ćemo formalno riješiti već na kraju ovog
poglavlja (odjeljak~\ref{sec:trikovi}).

\textbf{Povijest.} Ljudi teže „optimiranju” od davnina, ali korijeni
\textbf{moderne} optimizacije smještaju se u vrijeme \textbf{2.\ svjetskog
rata} --- tada nastaju metode \emph{operacijskih istraživanja} (engl.\
\emph{operations research}).

\textbf{Terminologija.} Pojam \textbf{„matematičko programiranje”} (engl.\
\emph{mathematical programming}) često se koristi kao sinonim za optimiranje.
Odatle nam dolaze i imena cijelih klasa problema koje ćemo obrađivati:

\begin{itemize}
  \item \emph{linear programming} (linearno programiranje, LP)
        $\to$ poglavlje~\ref{ch:lpqp},
  \item \emph{nonlinear programming} (nelinearno programiranje),
  \item \emph{quadratic programming} (kvadratno programiranje, QP)
        $\to$ poglavlje~\ref{ch:lpqp},
  \item \dots\ i tako dalje.
\end{itemize}

\begin{napomena}[riječ „programiranje”]
Riječ „programiranje” ovdje \textbf{nema veze s računalnim programiranjem}.
Potječe iz vojnog konteksta 1940-ih, gdje je „program” značio \emph{plan} ili
\emph{raspored} (npr.\ program opskrbe). To je uobičajen izvor zabune kod ljudi
koji se prvi put susreću s temom.
\end{napomena}

\textbf{Jezična opaska sa slajda} \str{7}\textbf{.} Predavač posebno ističe
riječ \emph{najoptimalnije} --- kao primjer \emph{pogrešne} upotrebe. Razlog:
„optimalno” već znači \emph{najbolje}. Ne postoji „optimalnije” ni
„najoptimalnije”, kao što ne postoji „najnajbolje”. Ispravno je reći
\textbf{optimalno}.

% =====================================================================
\section{Optimizacijski problem --- standardna forma}
\label{sec:standardna-forma}
\str{8--9}

Ovo je \textbf{središnja definicija cijelog kolegija}. Sve što slijedi u
poglavljima \ref{ch:definicije}--\ref{ch:robusno} samo je razrada ovog zapisa.

\subsection{Zapis s pojedinačnim ograničenjima}

\begin{definicija}[Optimizacijski problem]
\begin{equation}
\begin{aligned}
\minim & f(x) \\
\uzogr & h_i(x) = 0, \quad i = 1,\dots,m \\
       & g_j(x) \le 0, \quad j = 1,\dots,p
\end{aligned}
\end{equation}
\end{definicija}

\subsection{Objašnjenje svakog simbola od nule}

\textbf{Vektor projektnih varijabli $x$.}
\[
x = \begin{pmatrix} x_1 \\ x_2 \\ \vdots \\ x_n \end{pmatrix}
\]

\begin{itemize}
  \item To su \textbf{projektne varijable} (engl.\ \emph{decision variables},
        hrv.\ i \emph{varijable odluke}).
  \item Ovo su veličine koje \textbf{mi biramo} --- one su „nepoznanice”
        problema.
  \item $n$ je broj tih veličina. Vektor $x$ je dakle stupac od $n$ realnih
        brojeva; pišemo $x \in \R^n$.
  \item Notacija $\R$ označava skup realnih brojeva, a $\R^n$ skup svih
        uređenih $n$-torki realnih brojeva. Npr.\ $\R^2$ je obična ravnina.
\end{itemize}

\textbf{Funkcija cilja $f$.}
\[
f : \R^n \to \R
\]

\begin{itemize}
  \item Naziva se \textbf{funkcija cilja}, \textbf{ciljna funkcija} ili
        \textbf{kriterij} (engl.\ \emph{objective function}, \emph{cost
        function}).
  \item Zapis $f : \R^n \to \R$ čita se: „$f$ je funkcija koja uzima vektor iz
        $\R^n$ i vraća \textbf{jedan realan broj}”. Baš zato što vraća
        \textbf{jedan} broj, možemo dva različita izbora $x$ međusobno
        usporediti i reći koji je bolji.
  \item Konvencija je da \textbf{minimiziramo}. Ako želimo maksimizirati,
        koristimo trik iz odjeljka~\ref{sec:trikovi}.
\end{itemize}

\textbf{Ograničenja jednakosti $h_i$.}
\[
h_i : \R^n \to \R, \quad i = 1,\dots,m
\]

\begin{itemize}
  \item Engl.\ \emph{equality constraints}.
  \item Uvjet $h_i(x) = 0$ mora biti zadovoljen \textbf{točno}, bez ikakvog
        odstupanja.
  \item Tipično dolaze iz fizikalnih zakona ili definicijskih veza (npr.\
        „površina mora biti točno $A$”, „proizvodnja mora biti jednaka
        potrošnji”).
  \item $m$ je broj ograničenja jednakosti.
\end{itemize}

\textbf{Ograničenja nejednakosti $g_j$.}
\[
g_j : \R^n \to \R, \quad j = 1,\dots,p
\]

\begin{itemize}
  \item Engl.\ \emph{inequality constraints}.
  \item Uvjet $g_j(x) \le 0$ dopušta „rezervu”: dovoljno je da vrijednost ne
        prijeđe nulu.
  \item Tipično dolaze iz tehničkih ili resursnih ograničenja („naprezanje ne
        smije prijeći dopušteno”, „ne možemo proizvesti više od kapaciteta”).
  \item $p$ je broj ograničenja nejednakosti.
\end{itemize}

\textbf{Projektni prostor.} Općenito pišemo $x \in X$, gdje je $X$
\textbf{projektni prostor}. U svemu što slijedi u ovom kolegiju bit će
$X = \R^n$.

\subsection{Kompaktni („vektorski”) zapis}
\str{9}

Umjesto da nabrajamo ograničenja jedno po jedno, sve $h_i$ složimo u jedan
vektorski zapis, i isto za $g_j$:

\begin{equation}
\begin{aligned}
\minim & f(x) \\
\uzogr & h(x) = 0 \\
       & g(x) \le 0
\end{aligned}
\end{equation}

gdje je sada
\[
h : \R^n \to \R^m, \quad
h(x) = \begin{pmatrix} h_1(x) \\ \vdots \\ h_m(x) \end{pmatrix}, \qquad
g : \R^n \to \R^p, \quad
g(x) = \begin{pmatrix} g_1(x) \\ \vdots \\ g_p(x) \end{pmatrix}
\]

\begin{napomena}[kako čitati $g(x) \le 0$ kad je $g(x)$ vektor]
Nejednakost između vektora ovdje se čita \textbf{po komponentama}:
$g(x) \le 0$ znači $g_1(x) \le 0$ \emph{i} $g_2(x) \le 0$ \emph{i} \dots\
\emph{i} $g_p(x) \le 0$. Sve nejednakosti moraju vrijediti istovremeno. Nula na
desnoj strani je nul-vektor odgovarajuće dimenzije.
\end{napomena}

\subsection{Što je optimalno rješenje}
\str{9}

\begin{definicija}[Optimalno rješenje]
\textbf{Optimalno rješenje $x\opt$} je vektor projektnih varijabli koji, od
svih vektora iz projektnog prostora koji zadovoljavaju ograničenja, ima
\textbf{najmanju} vrijednost funkcije cilja $f$.
\end{definicija}

Uočite da definicija ima dva dijela i oba su nužna:

\begin{enumerate}
  \item $x\opt$ mora \textbf{sam zadovoljavati ograničenja} (inače nije
        kandidat),
  \item među svim takvim kandidatima ima \textbf{najmanji} $f$.
\end{enumerate}

Zvjezdica $\opt$ je standardna oznaka za optimalnu vrijednost; koristit ćemo je
kroz cijelu skriptu ($x\opt$, $f(x\opt)$, kasnije $p\opt$, $\lambda\opt$,
\dots).

\subsection{Dozvoljeni skup}
\str{78}

Skup svih $x$ koji zadovoljavaju \textbf{sva} ograničenja zove se
\textbf{dozvoljeni skup} (ili \emph{dozvoljeno područje}, \emph{dozvoljeni
projektni prostor}; engl.\ \emph{feasible set}):

\begin{definicija}[Dozvoljeni skup]
\begin{equation}
\Feas = \{\, x \mid h(x) = 0, \; g(x) \le 0 \,\}
\end{equation}
\end{definicija}

Notacija $\{\, x \mid \text{uvjet} \,\}$ čita se: „skup svih $x$ \textbf{takvih
da} vrijedi uvjet”. Uspravna crta $\mid$ znači „takvih da”.

Uz tu oznaku problem se piše najkraće:
\[
\min_{x \in \Feas} f(x)
\]

% =====================================================================
\section{Uvodni primjer s potpunim grafom}
\label{sec:uvodni-primjer}
\str{10--15}

Ovo je prvi konkretan problem u kolegiju i vrijedi ga proći vrlo pažljivo, jer
se na njemu vidi \emph{sve} iz prethodnog odjeljka odjednom.

\subsection{Postavka problema}
\str{10}

\begin{equation}
\begin{aligned}
\min_x \quad & (x_1 - 2)^2 + (x_2 - 1)^2 \\
\uzogr &
\begin{cases}
x_1^2 - x_2 \le 0 \\
x_1 + x_2 \le 2
\end{cases}
\end{aligned}
\end{equation}

\textbf{Projektni prostor} \str{10}: $X = \R^n$ s varijablama
$x = (x_1, x_2)$, dakle $n = 2$.

\subsection{Prepoznavanje elemenata standardne forme}
\str{11--12}

\textbf{Dozvoljeni skup} \str{11}:
\[
\Feas = \{\, x \mid g(x) \le 0 \,\}, \qquad
g(x) = \begin{pmatrix} g_1(x) \\ g_2(x) \end{pmatrix}
= \begin{pmatrix} x_1^2 - x_2 \\ x_1 + x_2 - 2 \end{pmatrix} \le 0
\]

Obratite pažnju na \textbf{prebacivanje na standardnu formu}: ograničenje je
zadano kao $x_1 + x_2 \le 2$, a standardna forma traži oblik $g_j(x) \le 0$.
Zato \textbf{oduzmemo 2 s obje strane}:
\[
x_1 + x_2 \le 2 \quad \Longleftrightarrow \quad x_1 + x_2 - 2 \le 0
\quad \Longrightarrow \quad g_2(x) = x_1 + x_2 - 2
\]

Prvo ograničenje već je u pravom obliku: $g_1(x) = x_1^2 - x_2$.

\textbf{Funkcija cilja} \str{12}:
\[
f(x) = (x_1 - 2)^2 + (x_2 - 1)^2
\]

\textbf{Geometrijsko značenje funkcije cilja.} $f(x)$ je \textbf{kvadrat
udaljenosti} točke $x = (x_1, x_2)$ od točke $(2, 1)$. (Po Pitagorinu poučku,
udaljenost je $\sqrt{(x_1-2)^2 + (x_2-1)^2}$, a $f$ je taj izraz na kvadrat.)
Dakle problem glasi: \textbf{„nađi točku dozvoljenog skupa koja je najbliža
točki $(2,1)$”}. Bez ograničenja rješenje bi trivijalno bilo $x = (2,1)$ s
$f = 0$; ograničenja to zabranjuju i zato problem nije trivijalan.

\subsection{Rješenje}
\str{15}

\[
x\opt = \begin{pmatrix} x_1\opt \\ x_2\opt \end{pmatrix}
= \begin{pmatrix} 1 \\ 1 \end{pmatrix}, \qquad f(x\opt) = 1
\]

\textbf{Provjera da je $x\opt$ doista dozvoljena točka} (uvrštavanje korak po
korak):
\begin{align}
g_1(x\opt) &= (x_1\opt)^2 - x_2\opt = 1^2 - 1 = 1 - 1 = 0 \le 0 \quad \checkmark\\
g_2(x\opt) &= x_1\opt + x_2\opt - 2 = 1 + 1 - 2 = 2 - 2 = 0 \le 0 \quad \checkmark
\end{align}

Oba ograničenja su zadovoljena \textbf{s jednakošću} --- kažemo da su
\textbf{aktivna} u $x\opt$ (taj ćemo pojam formalno uvesti u
poglavlju~\ref{ch:uvjeti}).

\textbf{Provjera vrijednosti funkcije cilja:}
\[
f(x\opt) = (1 - 2)^2 + (1 - 1)^2 = (-1)^2 + 0^2 = 1 + 0 = 1 \quad \checkmark
\]

\subsection{Graf}

Slajdovi 13, 14 i 15 grade upravo ovaj crtež u tri koraka: prvo ograničenje
$g_1$, zatim ograničenje $g_2$, pa oba zajedno s nivo krivuljama.

\begin{crtanje}[dozvoljeni skup i nivo krivulje]
\textbf{Korak 1 --- prazan koordinatni sustav.}
Nacrtajte vodoravnu os i označite je $x_1$, te okomitu os i označite je $x_2$.
Raspon kao na slajdu: $x_1$ od $-2$ do $3$, $x_2$ od $-2$ do $4$. Označite
jedinice na obje osi (cijeli brojevi). Koristite \textbf{jednako mjerilo} na
obje osi --- inače će kružnice iz koraka 5 izgledati kao elipse i izgubit će se
geometrijsko značenje.

\medskip
\textbf{Korak 2 --- prvo ograničenje $g_1$} (odgovara slajdu 13).
Granica ograničenja $g_1$ je krivulja gdje vrijedi jednakost:
\[
g_1(x) = 0 \quad \Longleftrightarrow \quad x_1^2 - x_2 = 0
\quad \Longleftrightarrow \quad x_2 = x_1^2
\]
To je \textbf{obična parabola} s tjemenom u ishodištu $(0,0)$, otvorena prema
gore. Nacrtajte je kroz nekoliko točaka koje lako izračunate:

\begin{center}
\begin{tabular}{lccccc}
\toprule
$x_1$ & $-2$ & $-1$ & $0$ & $1$ & $2$ \\
$x_2 = x_1^2$ & $4$ & $1$ & $0$ & $1$ & $4$ \\
\bottomrule
\end{tabular}
\end{center}

\emph{Koju stranu sjenčati?} Ograničenje traži $x_1^2 - x_2 \le 0$, tj.\
$x_2 \ge x_1^2$. Dakle dopušteno je ono što je \textbf{iznad} parabole.
Provjera testnom točkom: uzmimo $(0, 1)$ --- tada je $g_1 = 0^2 - 1 = -1 \le 0$
\checkmark, a točka $(0,1)$ leži iznad parabole. Dakle: \textbf{osjenčajte (kao
nedopušteno) područje ispod parabole}, a zadržite područje iznad nje.

\medskip
\textbf{Korak 3 --- drugo ograničenje $g_2$} (odgovara slajdu 14).
Granica je pravac:
\[
g_2(x) = 0 \quad \Longleftrightarrow \quad x_1 + x_2 - 2 = 0
\quad \Longleftrightarrow \quad x_2 = 2 - x_1
\]
Pravac s nagibom $-1$. Najlakše ga je nacrtati kroz dvije točke --- sjecišta s
osima:
\begin{itemize}
  \item za $x_1 = 0$: $x_2 = 2$ $\to$ točka $(0, 2)$,
  \item za $x_2 = 0$: $x_1 = 2$ $\to$ točka $(2, 0)$.
\end{itemize}
Povucite pravac kroz te dvije točke.

\emph{Koju stranu sjenčati?} Ograničenje traži $x_1 + x_2 - 2 \le 0$, tj.\
$x_2 \le 2 - x_1$, dakle dopušteno je \textbf{ispod} pravca. Provjera testnom
točkom: ishodište $(0,0)$ daje $g_2 = 0 + 0 - 2 = -2 \le 0$ \checkmark, a
ishodište je ispod pravca. Dakle: \textbf{osjenčajte (kao nedopušteno)
područje iznad pravca}.

\medskip
\textbf{Korak 4 --- odnos dviju krivulja i dozvoljeni skup.}
Parabola i pravac se sijeku ondje gdje vrijedi i $x_2 = x_1^2$ i
$x_2 = 2 - x_1$:
\[
x_1^2 = 2 - x_1 \;\Longrightarrow\; x_1^2 + x_1 - 2 = 0
\;\Longrightarrow\; (x_1 + 2)(x_1 - 1) = 0
\;\Longrightarrow\; x_1 = -2 \;\text{ ili }\; x_1 = 1
\]
Pripadni $x_2$:
\begin{itemize}
  \item za $x_1 = -2$: $x_2 = 2 - (-2) = 4$ $\to$ sjecište $(-2, 4)$,
  \item za $x_1 = 1$: $x_2 = 2 - 1 = 1$ $\to$ sjecište $(1, 1)$.
\end{itemize}
Dozvoljeni skup je presjek dvaju područja: \textbf{iznad parabole i ispod
pravca}. To je oblik \textbf{„leće”} (na slajdu bijelo područje) omeđen odozdo
parabolom, a odozgo pravcem, s vrhovima točno u dvjema sjecišnim točkama
$(-2,4)$ i $(1,1)$. Sve ostalo na crtežu je nedopušteno (na slajdu obojeno
crveno).

\medskip
\textbf{Korak 5 --- nivo krivulje funkcije cilja} (odgovara slajdu 15).
\textbf{Nivo krivulja} (engl.\ \emph{level curve}, \emph{contour}) funkcije $f$
za razinu $c$ je skup svih točaka u kojima $f$ ima \emph{istu} vrijednost $c$:
$\{\, x \mid f(x) = c \,\}$. (Pojam se detaljnije obrađuje u
poglavlju~\ref{ch:definicije}.) Ovdje:
\[
(x_1 - 2)^2 + (x_2 - 1)^2 = c
\]
To je \textbf{kružnica sa središtem u $(2,1)$ i polumjerom $\sqrt{c}$}.
Nacrtajte nekoliko koncentričnih kružnica oko točke $(2,1)$, npr.\ za
\[
c = 0{,}25 \;(r = 0{,}5), \quad c = 1 \;(r = 1), \quad c = 2 \;(r \approx 1{,}41),
\quad c = 4 \;(r = 2), \quad c = 9 \;(r = 3)
\]
Što je kružnica manja, to je vrijednost funkcije cilja manja --- jer je $f$
kvadrat udaljenosti od središta.

\medskip
\textbf{Korak 6 --- očitavanje optimuma.}
Označite točku $(2,1)$ (križić) --- to bi bio minimum \emph{da nema
ograničenja}, ali ona leži u crvenom (nedopuštenom) području: provjerimo,
$g_1(2,1) = 2^2 - 1 = 3 > 0$, dakle nedopustiva. Sada „napuhujte” kružnice od polumjera
$0$ prema van i gledajte \textbf{koja prva dotakne bijelu leću}. Ta prva dodirna
točka je optimum. To je točno kružnica $c = 1$, a dodiruje leću u vrhu $(1,1)$
--- dakle $x\opt = (1,1)$, $f(x\opt) = 1$.

\medskip
\textbf{Korak 7 --- što graf znači.}
Optimum leži u \textbf{vrhu} dozvoljenog skupa, ondje gdje se sijeku obje
granice, dakle \textbf{oba ograničenja su aktivna}. Optimalna vrijednost
$f(x\opt) = 1$ strogo je veća od nule (vrijednosti neograničenog minimuma) ---
to je „cijena” koju plaćamo zato što nas ograničenja tjeraju dalje od željene
točke $(2,1)$. U poglavlju~\ref{ch:uvjeti} naučit ćemo algebarski uvjet (KKT
uvjete) koji formalno opisuje upravo ovakvu situaciju.
\end{crtanje}

\begin{figure}[H]
\centering
\includegraphics[width=0.86\textwidth]{01-uvodni-primjer.png}
\caption{Dozvoljeni skup (bijela „leća”), nivo krivulje funkcije cilja
(koncentrične kružnice oko $(2,1)$) i optimum $x\opt=(1,1)$.
Nedopušteno područje je crveno.}
\end{figure}

\begin{napomena}[tiskarska pogreška na slajdu 15]
Na slajdu 15 prvo ograničenje je otipkano kao $g_1(x) = x_1^2 - x^2 \le 0$, dok
na slajdovima 10--14 svugdje piše $g_1(x) = x_1^2 - x_2 \le 0$. Ispravno je
$x_1^2 - x_2$ --- to potvrđuje i sam crtež na slajdu (parabola $x_2 = x_1^2$) i
rješenje $x\opt = (1,1)$.
\end{napomena}

% =====================================================================
\section{Gdje i zašto se optimira?}
\str{16--17}

\subsection{Optimiranje rada sustava / procesa („on-line”)}
\begin{itemize}
  \item Planiranje proizvodnje u elektroenergetskoj mreži
  \item Upravljanje motorom automobila
  \item Upravljanje tokovima snage i energijom u hibridnom automobilu
  \item Traženje najkraćeg puta (npr.\ GPS u autu, „preračunavam”)
  \item Gibanje robota prilikom obavljanja zadatka
  \item Plan bušenja rupa na ploči za elektroniku
  \item Web tražilica
  \item Optimalno upravljanje dinamičkim sustavima (bogata inženjerska /
        znanstvena grana za sebe)
  \item \dots
\end{itemize}

\subsection{U fundamentalnim fizikalnim zakonima / problemima}
\begin{itemize}
  \item Traženje \textbf{ravnotežnog položaja} minimizacijom potencijalne
        energije
  \item \textbf{Stabilnost dinamičkih sustava} --- traženje Ljapunovljeve
        funkcije
\end{itemize}

Ove dvije stavke nisu samo ilustracija: prva se detaljno rješava u
poglavlju~\ref{ch:uvjeti}, a druga je tema cijelog poglavlja~\ref{ch:dinamicki}.

\subsection{U prirodi}
\begin{itemize}
  \item Oblik stabla (optimiranje oblika obzirom na vlastitu težinu i vanjska
        opterećenja)
  \item Struktura kosti
  \item Trajektorija šišmiša u lovu na plijen
  \item Oblik jabuke
  \item Prilagodba vrste na uvjete života u zadanoj okolini
  \item \dots
\end{itemize}

\subsection{Optimalno projektiranje}
\begin{itemize}
  \item Optimalna raspodjela vjetroturbina u vjetroparku
  \item Optimiranje topologije, oblika, dimenzija mehanizama i nosivih
        konstrukcija --- \textbf{u ovom kolegiju}
  \item \dots
\end{itemize}

\dots\ i još puno, puno toga.

% =====================================================================
\section{Sadržaj kolegija i ishodi učenja}
\str{18--32}

Kompletan popis tema sa slajdova (gradi se inkrementalno kroz 15 slajdova):

\begin{enumerate}
  \item \textbf{Uvod}: Ciljevi kolegija i ishodi učenja.
  \item \textbf{Optimizacijski problem}: definicije i klasifikacija. Konveksni
        opt.\ problemi.
  \item \textbf{Uvjeti optimalnosti.} Geometrijska interpretacija i primjene.
  \item \textbf{Lagrangeova dualnost}, alternative i analiza osjetljivosti.
  \item \textbf{Linearno programiranje.} Pristupi rješavanju i primjeri.
  \item \textbf{Kvadratno programiranje} s primjenama. Formulacija i značajke.
  \item \textbf{Linearne matrične nejednadžbe i semidefinitno programiranje.}
  \item \textbf{Numerički algoritmi} za rješavanje opt.\ problema (pregled).
  \item \textbf{Optimiranje mehatroničkih sustava} --- od linearnog do
        semidefinitnog programiranja: primjeri iz prakse.
  \item \textbf{Formulacija stabilnosti i disipativnost dinamičkih sustava} kao
        opt.\ problema.
  \item Formulacije kriterija i analiza učinkovitosti dinamičkih sustava.
  \item Formulacije \textbf{optimalnog upravljanja} i pristupi rješavanju.
  \item \textbf{Robusno optimiranje 1}
  \item \textbf{Robusno optimiranje 2}
  \item \textbf{Višeciljna optimizacija} i višeciljno upravljanje
\end{enumerate}

\textbf{Kako se to preslikava na ovu skriptu:} teme 1--2 $\to$
poglavlja~\ref{ch:uvod}--\ref{ch:definicije}; tema 3 $\to$
poglavlje~\ref{ch:uvjeti}; tema 4 $\to$ poglavlje~\ref{ch:dualnost}; teme 5--6
$\to$ poglavlje~\ref{ch:lpqp}; teme 7 i 9 $\to$ poglavlje~\ref{ch:qcpsdp}; teme
10--12 $\to$ poglavlje~\ref{ch:dinamicki}; teme 13--14 $\to$
poglavlje~\ref{ch:robusno}.

% =====================================================================
\section{Motivacijski primjeri iz prakse}
\str{33--43}

\subsection{Planiranje proizvodnje u elektroenergetskoj mreži}
\str{33--34}

\textbf{Problem:} planiranje rada električnih generatora.

Sa slajda, sa svrstavanjem svake stavke u element standardne forme:

\begin{center}
\begin{tabular}{p{7.4cm}p{6.0cm}}
\toprule
\textbf{Sa slajda} & \textbf{Element opt.\ problema} \\
\midrule
tipično noć prije dana proizvodnje & (vremenski okvir odluke) \\
proizvodnja $=$ predviđenoj potrošnji & \textbf{ograničenje jednakosti} $h(x)=0$ \\
minimiziranje troškova proizvodnje & \textbf{funkcija cilja} $f(x)$ \\
ograničenja na tok snage u dalekovodima & \textbf{ograničenja nejednakosti} $g(x)\le 0$ \\
planovi u slučaju kvara & \textbf{robusnost} (poglavlje~\ref{ch:robusno}) \\
\bottomrule
\end{tabular}
\end{center}

Dodatno \str{34}:
\begin{itemize}
  \item Optimiranje (optimalno upravljanje) je \textbf{„skrivena tehnologija”}
        za razvoj naprednih mreža u budućnosti.
  \item Matematička formulacija tržišta i kreiranja tržišnih cijena:
        \textbf{Lagrangeova dualnost}.
\end{itemize}

Zadnja rečenica je najava Vježbi 3, gdje se cijeli „Primjer 1: Tržište i
formiranje cijena” rješava upravo preko dualnosti --- vidi
poglavlje~\ref{ch:dualnost}.

\subsection{Upravljanje tokovima snage i energijom u hibridnom automobilu}
\str{35}

\begin{itemize}
  \item Nije dobro imati napunjenu bateriju: želimo je (polu)praznu da bismo je
        mogli puniti prilikom kočenja (rekuperacija).
  \item \textbf{dobro $\to$ optimalno}
  \item Optimiranje je „skrivena tehnologija” koja daje „inteligenciju”
        procesu: pojmovi \emph{dobro--loše} moraju biti vezani uz neki kriterij
        $\to$ \textbf{formalizacija funkcije cilja}.
\end{itemize}

\subsection{Prediktivno upravljanje dinamičkim sustavima}
\str{36--37}

Slajdovi su ilustracija koncepta MPC-a (engl.\ \emph{Model Predictive
Control}): u svakom koraku se, na temelju modela, riješi optimizacijski problem
na konačnom horizontu predikcije, primijeni se prvi upravljački potez, pa se
postupak ponovi.

\subsection{Optimalno upravljanje --- pozicioniranje glave tvrdog diska}
\str{38}

Slajd prikazuje mehanizam tvrdog diska (ruka s glavom, \emph{pivot},
\emph{voice coil motor} VCM) i regulacijsku petlju: regulator $K$ prima signal
pogreške $e$ i daje upravljanje $u$; $y_r$ je referenca, $y$ izlaz, $d$
poremećaj.

Podaci sa slajda: ${>}15\,000$ o/min, $12\,000$ traka na $1$ cm.

\textbf{Formulacija problema (sa slajda):}
\begin{equation}
\min_K \; \max_{\|d\| \ne 0} \frac{\|e\|}{\|d\|}, \qquad
K \in \{\,\text{stabilizirajući regulatori}\,\}
\end{equation}

\textbf{Kako to čitati.} $\|\cdot\|$ je \textbf{norma} --- mjera „veličine”
signala (formalno u poglavlju~\ref{ch:definicije}). Omjer $\|e\|/\|d\|$ je
„koliko se poremećaj $d$ pojača u pogrešku $e$”. Unutarnji $\max$ uzima
\textbf{najgori mogući} poremećaj, a vanjski $\min$ bira regulator koji taj
najgori slučaj čini što manjim. Ovakav \textbf{min--max} zapis je tipičan za
robusno projektiranje i vratit ćemo mu se u
poglavljima~\ref{ch:dinamicki}~i~\ref{ch:robusno}.

\subsection{Ravnotežni položaj sustava u potencijalnom polju}
\str{39--40}

\textbf{Postavka} \str{39}\textbf{.} Točka mase $m$ obješena je u vertikalnoj
ravnini o dvije elastične opruge zanemarive mase s konstantama $k_1$, $k_2$ i
duljinama $l_1$, $l_2$, dok je razmak objesišta jednak $d$.

Geometrija sa slajda: koordinatni sustav $(x, y)$; lijevo objesište je u
ishodištu $(0,0)$, desno objesište je na osi $x$ na udaljenosti $d$, tj.\ u
$(d, 0)$; masa $m$ visi ispod, u točki $(x, y)$ (s $y < 0$). Lijeva opruga ima
parametre $k_1, l_1$, desna $k_2, l_2$.

\textbf{Zadatak:} izračunati ravnotežni položaj mase.

\textbf{Rješenje} \str{40}\textbf{.} Potencijalna energija mase $m$ u položaju
$(x,y)$ je
\begin{equation}
V(x,y) = mgy
+ \frac{k_1}{2}\left(l_1 - \sqrt{x^2 + y^2}\right)^2
+ \frac{k_2}{2}\left(l_2 - \sqrt{(d-x)^2 + y^2}\right)^2
\end{equation}

\textbf{Rastav izraza član po član} (svaki je član oblika „energija”):

\begin{itemize}
  \item $mgy$ --- \textbf{gravitacijska potencijalna energija} mase u polju sile
        teže ($g$ = ubrzanje sile teže). Budući da masa visi ispod osi $x$, $y$
        je negativan pa je taj član negativan --- spuštanjem mase energija pada.
  \item $\sqrt{x^2 + y^2}$ --- trenutna \textbf{duljina lijeve opruge}, tj.\
        udaljenost od objesišta $(0,0)$ do mase $(x,y)$ (Pitagorin poučak).
  \item $l_1 - \sqrt{x^2+y^2}$ --- \textbf{deformacija} lijeve opruge (razlika
        nedeformirane duljine $l_1$ i trenutne duljine).
  \item $\frac{k_1}{2}(\cdot)^2$ --- \textbf{elastična energija} lijeve opruge;
        standardna formula $E = \tfrac{1}{2}k\,\Delta l^2$.
  \item Analogno za desnu oprugu, samo je njezino objesište u $(d, 0)$, pa je
        udaljenost do mase $\sqrt{(d-x)^2 + y^2}$.
\end{itemize}

\textbf{Ravnotežna točka} $(x\opt, y\opt)$ je ona koja \textbf{minimizira}
potencijalnu energiju:
\begin{equation}
(x\opt, y\opt) = \argmin_{x,\,y} \; V(x,y)
\end{equation}

\begin{napomena}[notacija $\argmin$]
$\min_x V(x)$ je \textbf{najmanja vrijednost} koju $V$ poprima (broj), dok je
$\argmin_x V(x)$ \textbf{argument} u kojem se taj minimum postiže (točka). Za
nas je obično zanimljiv $\argmin$ --- želimo znati \emph{gdje} je optimum, a ne
samo \emph{koliko} iznosi.
\end{napomena}

Ovo je važan konceptualni primjer: \textbf{fizikalni zakon (ravnoteža) izražen
je kao optimizacijski problem bez ograničenja.} Isti primjer se u
poglavlju~\ref{ch:uvjeti} proširuje na slučaj s ograničenjima i rješava preko
uvjeta optimalnosti.

\subsection{Primjeri iz prirode}
\str{41}

Slajd je fotografska ilustracija tvrdnji iz odjeljka o prirodi (oblici koje
priroda „optimira”). Nema dodatnih formula.

\subsection{Raspodjela vjetroturbina u vjetroparku}
\str{42--43}

Slajd 43 prikazuje dva 3D grafa (površina prinosa nad tlocrtom parka, s
ucrtanim pozicijama turbina kao točkama i označenim smjerom sjevera \textbf{N}):

\begin{itemize}
  \item lijevo: \textbf{optimalni položaji sa zanemarenim interakcijama preko
        vrtloga},
  \item desno: \textbf{optimalni položaji s uračunatim interakcijama preko
        vrtloga}.
\end{itemize}

\textbf{Poanta:} ista funkcija cilja (ukupna proizvedena energija), ali bogatiji
model (uzimanje u obzir da turbina u vjetrenoj sjeni druge turbine daje manje)
daje \textbf{bitno drukčiji optimalni raspored}. Model ulazi u optimizaciju
kroz ograničenja i funkciju cilja --- kvaliteta modela izravno određuje
kvalitetu optimuma.

% =====================================================================
\section{Primjer: problem proizvodnje i distribucije robe}
\str{44--45}

Ovo je prvi potpuno formuliran „inženjersko-poslovni” problem u kolegiju i
ujedno prvi primjer \textbf{linearnog programa} (LP), koji ćemo formalno
obraditi u poglavlju~\ref{ch:lpqp}.

\subsection{Opis problema}
\str{44}

\begin{itemize}
  \item Tvrtka ima \textbf{2 tvornice} i \textbf{7 trgovina}.
  \item U jednom proizvodnom vremenskom intervalu tvornica $i$ može proizvesti
        do $a_i$ tona.
  \item Svaka trgovina $j$ redovito ima potražnju od $b_j$ tona.
  \item Troškovi transporta od tvornice $i$ do trgovine $j$ su $c_{ij}$
        kuna/toni.
  \item \textbf{Kako optimalno voditi tvrtku (minimizirati trošak)?}
\end{itemize}

\subsection{Projektne varijable}
\str{45}

\[
x_{ij} \in \R: \text{ koliko tvornica } i \text{ dobavlja robe trgovini } j,
\qquad x_{ij} \ge 0
\]

Budući da imamo $2$ tvornice i $7$ trgovina, ukupno je $2 \times 7 = 14$
varijabli. Uvjet $x_{ij} \ge 0$ je fizikalne prirode: ne može se isporučiti
negativna količina.

\subsection{Potpuna formulacija}
\str{45}

\begin{equation}
\begin{aligned}
\min \quad & \sum_{i=1}^{2} \sum_{j=1}^{7} c_{ij} x_{ij} \\
\uzogr & \sum_{j=1}^{7} x_{ij} \le a_i, \qquad i = 1, 2 \\
       & \sum_{i=1}^{2} x_{ij} \ge b_j, \qquad j = 1, \dots, 7 \\
       & x_{ij} \ge 0
\end{aligned}
\end{equation}

\textbf{Čitanje svakog retka:}

\begin{itemize}
  \item \textbf{Funkcija cilja} $\sum_i \sum_j c_{ij} x_{ij}$: za svaku moguću
        relaciju tvornica $\to$ trgovina pomnožimo prevezenu količinu $x_{ij}$
        [t] s jediničnim troškom $c_{ij}$ [kn/t] i sve zbrojimo. Rezultat je
        \textbf{ukupan trošak transporta} u kunama. (Simbol $\sum$ znači
        „zbroji po svim vrijednostima indeksa”; dvostruka suma znači da
        prolazimo po svim parovima $(i,j)$, dakle po svih 14 relacija.)
  \item \textbf{Kapacitet tvornice} $\sum_{j} x_{ij} \le a_i$: sve što tvornica
        $i$ pošalje (zbroj po svim trgovinama) ne smije prijeći njezin kapacitet
        $a_i$. Dva takva ograničenja (po jedno za svaku tvornicu).
  \item \textbf{Potražnja trgovine} $\sum_{i} x_{ij} \ge b_j$: sve što trgovina
        $j$ primi (zbroj po objema tvornicama) mora biti barem $b_j$. Sedam
        takvih ograničenja.
  \item \textbf{Nenegativnost} $x_{ij} \ge 0$: 14 ograničenja.
\end{itemize}

\begin{napomena}[prevođenje u standardnu formu]
Standardna forma traži $g_j(x)\le 0$, pa nejednakost „$\ge$” pomnožimo s $-1$
(što okreće smjer): $\sum_i x_{ij} \ge b_j \Leftrightarrow
b_j - \sum_i x_{ij} \le 0$, i $x_{ij} \ge 0 \Leftrightarrow -x_{ij} \le 0$.
Kapacitet je već u pravom obliku nakon premještanja:
$\sum_j x_{ij} - a_i \le 0$.
\end{napomena}

Uočite: i funkcija cilja i sva ograničenja su \textbf{linearne} funkcije
varijabli $x_{ij}$. To je definicija linearnog programa
(poglavlje~\ref{ch:lpqp}).

% =====================================================================
\section{Klasični primjeri optimiranja s potpunim rješenjima}
\str{46--52}

Slajdovi 46--52 rješavaju šest školskih ekstremalnih problema. Svi se rješavaju
istim receptom, pa ga vrijedi jednom izreći:

\begin{definicija}[Recept klasične analize --- ekstrem funkcije jedne varijable]
\begin{enumerate}
  \item Uvedi varijablu i izrazi \textbf{funkciju cilja} preko nje.
  \item Ako postoji ograničenje jednakosti, iz njega \textbf{izrazi jednu
        varijablu} i \textbf{uvrsti} u funkciju cilja (eliminacija) --- tako
        problem s ograničenjem postaje problem bez ograničenja.
  \item Deriviraj i izjednači s nulom: $\dfrac{\dd F}{\dd x} = 0$ $\to$
        \textbf{stacionarne točke}.
  \item Provjeri drugom derivacijom: $F'' > 0$ $\to$ \textbf{minimum},
        $F'' < 0$ $\to$ \textbf{maksimum}.
  \item Izračunaj vrijednost cilja u nađenoj točki.
\end{enumerate}
\end{definicija}

Ovaj recept vrijedi za jednostavne, glatke probleme jedne varijable. U
poglavlju~\ref{ch:uvjeti} generalizirat ćemo ga na više varijabli i na probleme
gdje se ograničenja \emph{ne mogu} tako lako eliminirati --- tada nastupaju
Lagrangeovi multiplikatori i KKT uvjeti.

\subsection{Optimalno iskorištavanje materijala --- posuda iz limene ploče}
\str{46--47}

\textbf{Problem.} Iz kvadratne limene ploče stranice $a$ potrebno je izrezati
uglove tako da se iz preostalog dijela načini posuda što većeg volumena.
Izrezuju se četiri kvadratića stranice $x$ u uglovima; preostali „križ” se
savine prema gore.

\textbf{Ciljna funkcija je volumen:}
\[
V = (a - 2x)^2 \cdot x
\]

\emph{Odakle to:} nakon izrezivanja uglova stranice $x$ i savijanja, baza posude
je kvadrat stranice $a - 2x$ (s obje strane oduzimamo po $x$), a visina posude
je $x$.

\textbf{Derivacija:}
\[
\frac{\dd V}{\dd x} = 12x^2 - 8ax + a^2
\]

\emph{Provjera računa} (raspis međukoraka koji na slajdu nije napisan):
\begin{align}
V &= x(a-2x)^2 = x(a^2 - 4ax + 4x^2) = a^2 x - 4a x^2 + 4x^3 \\
\frac{\dd V}{\dd x} &= a^2 - 8ax + 12x^2 \quad \checkmark
\end{align}

\textbf{Izjednačavanje s nulom} daje kvadratnu jednadžbu
$12x^2 - 8ax + a^2 = 0$, čija su rješenja (sa slajda):
\[
x_{1,2} = \frac{a}{3} \pm \frac{a}{6}; \qquad x_1 = a/2, \quad x_2 = a/6
\]

\emph{Provjera preko formule za kvadratnu jednadžbu:}
\begin{align}
x_{1,2} &= \frac{8a \pm \sqrt{(8a)^2 - 4\cdot 12 \cdot a^2}}{2 \cdot 12}
= \frac{8a \pm \sqrt{64a^2 - 48a^2}}{24}
= \frac{8a \pm \sqrt{16a^2}}{24}
= \frac{8a \pm 4a}{24} \\
x_1 &= \frac{12a}{24} = \frac{a}{2}, \qquad
x_2 = \frac{4a}{24} = \frac{a}{6} \quad \checkmark
\end{align}

\textbf{Provjera drugom derivacijom:}
\[
V'' = \frac{\dd^2 V}{\dd x^2} = 24x - 8a
\]
\begin{align}
V''\!\left(\tfrac{a}{2}\right) &= 24 \cdot \tfrac{a}{2} - 8a = 12a - 8a = 4a > 0
\;\to\; \textbf{min.} \\
V''\!\left(\tfrac{a}{6}\right) &= 24 \cdot \tfrac{a}{6} - 8a = 4a - 8a = -4a < 0
\;\to\; \textbf{max.}
\end{align}

\textbf{Vrijednosti:}
\[
V_{\max} = \left(a - \frac{a}{3}\right)^{2} \cdot \frac{a}{6} = \frac{2}{27}a^3;
\qquad V_{\min} = 0
\]

\emph{Provjera:} za $x = a/6$ je $a - 2x = a - a/3 = \tfrac{2a}{3}$, pa je
$V = \left(\tfrac{2a}{3}\right)^2 \cdot \tfrac{a}{6}
= \tfrac{4a^2}{9}\cdot\tfrac{a}{6} = \tfrac{4a^3}{54} = \tfrac{2a^3}{27}$
\checkmark. Za $x = a/2$ je $a - 2x = 0$, pa je $V = 0$ --- geometrijski jasno:
izrezali smo cijelu ploču, baza je nestala.

\subsection{Minimizacija troškova izgradnje --- tlocrt kuće}
\str{48}

\textbf{Problem.} Kuća ima 3 prostorije površine $A$ [m\textsuperscript{2}].
Troškovi izgradnje vanjskih zidova iznose $C$ [EUR/m], a troškovi za
unutrašnje zidove iznose $2/3$ tih troškova. Odrediti mjere tlocrta tako da
troškovi izgradnje zidova budu minimalni.

Geometrija sa slajda: pravokutnik širine $x$ i visine $y$; unutar njega jedan
okomiti zid (duljine $y$) koji odvaja desni dio, te u lijevom dijelu vodoravni
zid duljine $2x/5$ --- tako se dobiju 3 prostorije.

\textbf{Ograničenje jednakosti:}
\[
A = x\,y \; [\text{m}^2] \quad \to \quad y = A/x
\]

\textbf{Ciljna funkcija $\to$ minimalni troškovi:}
\[
F(x) = C\left(2x + 2\frac{A}{x}\right)
     + \frac{2}{3}C\left(\frac{2}{5}x + \frac{A}{x}\right)
\]

\emph{Rastav:} prvi član su vanjski zidovi --- opseg pravokutnika je
$2x + 2y = 2x + 2A/x$, uz cijenu $C$. Drugi član su unutarnji zidovi ---
vodoravni duljine $\tfrac{2}{5}x$ i okomiti duljine $y = A/x$, uz cijenu
$\tfrac{2}{3}C$.

\textbf{Sređivanje (sa slajda):}
\[
F(x) = \frac{2}{3}C\left(\frac{17}{5}x + 4\frac{A}{x}\right)
\]

\emph{Raspis međukoraka koji na slajdu nije prikazan:}
\begin{align}
F(x) &= 2Cx + \frac{2CA}{x} + \frac{2}{3}C \cdot \frac{2}{5}x
        + \frac{2}{3}C\frac{A}{x} \\
     &= C\underbrace{\left(2 + \frac{4}{15}\right)}_{=\frac{30+4}{15}=\frac{34}{15}}x
      + C\underbrace{\left(2 + \frac{2}{3}\right)}_{=\frac{6+2}{3}=\frac{8}{3}}\frac{A}{x}
      = \frac{34}{15}Cx + \frac{8}{3}C\frac{A}{x}
\end{align}
a s druge strane
\[
\frac{2}{3}C\left(\frac{17}{5}x + \frac{4A}{x}\right)
= \frac{2 \cdot 17}{3 \cdot 5}Cx + \frac{2\cdot 4}{3}C\frac{A}{x}
= \frac{34}{15}Cx + \frac{8}{3}C\frac{A}{x} \quad \checkmark
\]

\textbf{Konstante prebacimo na lijevu stranu} (jer konstantni pozitivni faktor
ne mijenja položaj minimuma --- ovo je važan i često korišten trik):
\[
F_0 = \frac{3F(x)}{2C} = \left(\frac{17}{5}x + 4\frac{A}{x}\right)
\]

\textbf{Derivacija i izjednačavanje s nulom:}
\[
\frac{\dd F_0}{\dd x} = \frac{17}{5} - \frac{4A}{x^2} = 0
\quad \to \quad x^2 = \frac{20}{17}A = \frac{20}{17}xy
\]

\emph{Međukorak:} iz $\tfrac{17}{5} = \tfrac{4A}{x^2}$ slijedi
$x^2 = \tfrac{4A \cdot 5}{17} = \tfrac{20A}{17}$; zatim se uvrsti $A = xy$.

\textbf{Odakle treba biti:}
\[
\frac{x}{y} = \frac{20}{17}
\]

\emph{Međukorak:} iz $x^2 = \tfrac{20}{17}xy$ podijelimo obje strane s $x$
(dopušteno jer $x > 0$) i dobijemo $x = \tfrac{20}{17}y$, tj.\ $x/y = 20/17$.

\subsection{Maksimizacija volumena lijevka I}
\str{49}

\textbf{Problem.} Neka se iz okrugle ploče izreže isječak tako da se iz ostatka
načini lijevak (stožac) što većeg volumena.

Geometrija: iz kruga polumjera $r$ izreže se kružni isječak, a preostali dio
(središnjeg kuta $\varphi$) se smota u plašt stošca. Nastali stožac ima polumjer
baze $R$, izvodnicu $r$ (jer je to polumjer polazne ploče) i visinu $h$.

\textbf{Polumjer baze (1):} duljina luka preostalog dijela postaje opseg baze:
\begin{equation}
2R\pi = r\varphi \quad \to \quad R = \frac{r\varphi}{2\pi}
\end{equation}

\textbf{Visina stošca (2):} iz pravokutnog trokuta (izvodnica, polumjer baze,
visina):
\begin{equation}
h = \sqrt{r^2 - R^2}
\end{equation}

\textbf{Ciljna funkcija (3):} volumen stošca
\begin{equation}
V = \frac{\pi R^2 h}{3}
\end{equation}

\textbf{Uvrštenjem (1) i (2) u (3):}
\[
V = \frac{r^3}{12\pi}\varphi^2 \sqrt{1 - \frac{\varphi^2}{4\pi^2}}
\]

\emph{Raspis međukoraka:}
\begin{align}
R^2 &= \frac{r^2\varphi^2}{4\pi^2}, \qquad
h = \sqrt{r^2 - \frac{r^2\varphi^2}{4\pi^2}}
  = r\sqrt{1 - \frac{\varphi^2}{4\pi^2}} \\
V &= \frac{\pi}{3}\cdot\frac{r^2\varphi^2}{4\pi^2}\cdot
     r\sqrt{1 - \frac{\varphi^2}{4\pi^2}}
   = \frac{r^3\varphi^2}{12\pi}\sqrt{1 - \frac{\varphi^2}{4\pi^2}}
   \quad \checkmark
\end{align}

\textbf{Prebacivanjem konstanti na lijevu stranu i kvadriranjem, uređena ciljna
funkcija je:}
\[
F(\varphi) = \left(\frac{12\pi V}{r^3}\right)^{2}
= \varphi^4\left(1 - \frac{\varphi^2}{4\pi^2}\right)
\]

\begin{napomena}[zašto smijemo kvadrirati]
Funkcija $t \mapsto t^2$ je \textbf{strogo rastuća} za $t \ge 0$, a $V \ge 0$.
Strogo rastuća transformacija ne mijenja \emph{položaj} maksimuma, samo njegovu
vrijednost. Isto vrijedi za množenje pozitivnom konstantom. Ovaj postupak
(uklanjanje korijena kvadriranjem) koristit ćemo često --- vidi i Vježbe 2, gdje
se isto radi s normama.
\end{napomena}

\textbf{Derivacija:}
\[
\frac{\dd F}{\dd \varphi} = 4\varphi^3 - \frac{3}{2\pi^2}\varphi^5 = 0
\]

\emph{Provjera:} $F(\varphi) = \varphi^4 - \dfrac{\varphi^6}{4\pi^2}$, pa je
$F' = 4\varphi^3 - \dfrac{6\varphi^5}{4\pi^2}
    = 4\varphi^3 - \dfrac{3\varphi^5}{2\pi^2}$ \checkmark

\textbf{Iz toga slijedi:}
\[
\varphi^2 = \frac{8}{3}\pi^2 \quad \to \quad
\varphi = \sqrt{\frac{2}{3}}\cdot 2\pi \,\text{rad} = 293{,}939^\circ
\]

\emph{Međukoraci:} podijelimo jednadžbu s $\varphi^3$ (za $\varphi \ne 0$):
$4 = \tfrac{3}{2\pi^2}\varphi^2$, dakle $\varphi^2 = \tfrac{8\pi^2}{3}$. Zatim
$\varphi = \sqrt{\tfrac{8\pi^2}{3}} = 2\pi\sqrt{\tfrac{2}{3}}$. Brojčano:
$\sqrt{2/3} = 0{,}8165$, pa $\varphi = 0{,}8165 \cdot 6{,}2832 = 5{,}1302$ rad;
u stupnjevima $5{,}1302 \cdot \tfrac{180}{\pi} = 293{,}94^\circ$ \checkmark

\subsection{Maksimizacija volumena lijevka II}
\str{50}

\textbf{Problem.} Kakav oblik treba imati lijevak (stožac) da kod \textbf{zadane
površine plašta} ima što veći volumen?

\textbf{Površina plašta:}
\[
A = 2\pi r \cdot R / 2 = \pi R r \quad \to \quad r = \frac{A}{\pi R}
\]

(Ovdje je $r$ izvodnica stošca, $R$ polumjer baze.)

\textbf{Volumen:}
\[
V = \frac{\pi R^2}{3}\sqrt{r^2 - R^2}
  = \frac{\pi R^2}{3}\sqrt{\frac{A^2}{\pi^2R^2} - R^2}, \qquad
V = \frac{R}{3}\sqrt{A^2 - \pi^2R^4}
\]

\emph{Raspis zadnjeg koraka:} izlučimo $\tfrac{1}{\pi^2R^2}$ iz korijena:
\begin{align}
\sqrt{\frac{A^2}{\pi^2R^2} - R^2} &= \frac{1}{\pi R}\sqrt{A^2 - \pi^2 R^4} \\
V &= \frac{\pi R^2}{3}\cdot\frac{1}{\pi R}\sqrt{A^2 - \pi^2R^4}
   = \frac{R}{3}\sqrt{A^2 - \pi^2R^4} \quad \checkmark
\end{align}

\textbf{Volumen $V$ bit će maksimalan kad i $(3V)^2$}, pa je (isti trik kao
prije):
\[
F(R) = (3V)^2 = A^2R^2 - \pi^2R^6
\]

\textbf{Derivacija:}
\[
\frac{\dd F}{\dd R} = 2A^2R - 6\pi^2R^5 = 0 \quad \to \quad
A = \sqrt{3\pi^2R^4} = \pi\sqrt{3}\,R^2
\]

\emph{Međukorak:} podijelimo s $2R$ ($R \ne 0$): $A^2 = 3\pi^2R^4$, pa
korjenovanjem $A = \pi R^2\sqrt{3}$.

\textbf{Iz $R^2 = \dfrac{A}{\pi\sqrt{3}}$ i
$h^2 = r^2 - R^2 = \dfrac{A^2}{\pi^2R^2} - R^2$, slijedi:}
\[
\frac{h^2}{R^2} = 2 \quad \to \quad \frac{h}{R} = \sqrt{2}
\]

\emph{Raspis:}
\[
\frac{h^2}{R^2} = \frac{1}{R^2}\left(\frac{A^2}{\pi^2R^2} - R^2\right)
= \frac{A^2}{\pi^2R^4} - 1
\overset{A^2 = 3\pi^2R^4}{=} \frac{3\pi^2R^4}{\pi^2R^4} - 1 = 3 - 1 = 2
\quad \checkmark
\]

\textbf{Zaključak:} optimalan lijevak ima omjer visine i polumjera baze
$h/R = \sqrt{2}$.

\subsection{Dimenzioniranje otvora vrata}
\str{51}

\textbf{Problem.} Kako visoka moraju biti vrata tornja ($h$) da se ljestve
duljine $L$ mogu uvući u toranj?

Geometrija sa slajda: zid tornja je okomit; ljestve $\overline{AB} = L$
naslonjene su tako da im donji kraj $B$ klizi po podu na udaljenosti $x$ od
zida, a gornji kraj $A$ je uz vanjski zid; $a$ je vodoravni odmak zida, $y$ je
visina na kojoj ljestve prelaze preko brida vrata, a $h$ je tražena visina
vrata.

\textbf{Rješenje:} traži se \textbf{maksimalna visina $y$ za vrijeme uvlačenja}
ljestvi u toranj.

\textbf{Iz sličnosti trokuta (1):}
\begin{equation}
\frac{y}{x-a} = \frac{\sqrt{L^2 - x^2}}{x}
\quad \to \quad
y = \frac{x-a}{x}\sqrt{L^2 - x^2}
\end{equation}

(Ovdje je $\sqrt{L^2 - x^2}$ visina gornjeg kraja ljestvi $A$ nad podom ---
Pitagorin poučak na trokutu s hipotenuzom $L$ i vodoravnom katetom $x$.)

\textbf{Derivacija i stacionarna točka (2):}
\begin{equation}
\frac{\dd y}{\dd x} = \frac{aL^2 - x^3}{x^2\sqrt{L^2 - x^2}} = 0
\quad \to \quad x_{\mathrm m} = \sqrt[3]{aL^2}
\end{equation}

\emph{Međukorak:} razlomak je nula kad je brojnik nula: $aL^2 - x^3 = 0
\Rightarrow x^3 = aL^2 \Rightarrow x = \sqrt[3]{aL^2}$.

\textbf{Uvrštenjem (2) u (1):}
\[
h \ge y_{\mathrm m} = \left[1 - \left(\frac{a}{L}\right)^{2/3}\right]^{3/2}\cdot L
\]

\emph{Raspis uvrštavanja (na slajdu nije prikazan):} uz
$x_{\mathrm m} = a^{1/3}L^{2/3}$,
\begin{align}
\frac{a}{x_{\mathrm m}} &= \frac{a}{a^{1/3}L^{2/3}} = \frac{a^{2/3}}{L^{2/3}}
= \left(\frac{a}{L}\right)^{2/3} \\
\sqrt{L^2 - x_{\mathrm m}^2} &= \sqrt{L^2 - a^{2/3}L^{4/3}}
= L\sqrt{1 - \left(\frac{a}{L}\right)^{2/3}} \\
y_{\mathrm m} &= \left(1 - \frac{a}{x_{\mathrm m}}\right)
                 \sqrt{L^2 - x_{\mathrm m}^2}
= \left[1 - \left(\tfrac{a}{L}\right)^{2/3}\right]
  \cdot L\left[1 - \left(\tfrac{a}{L}\right)^{2/3}\right]^{1/2}
= L\left[1 - \left(\tfrac{a}{L}\right)^{2/3}\right]^{3/2} \; \checkmark
\end{align}

\textbf{Brojčani primjer sa slajda:}
\[
a = 2{,}7\,\text{m}; \quad L = 12{,}5\,\text{m}
\;\to\; y_{\mathrm m} = 6{,}4\,\text{m}; \quad x_{\mathrm m} = 7{,}5\,\text{m}
\]

\emph{Provjera računa:}
\begin{align}
\frac{a}{L} &= \frac{2{,}7}{12{,}5} = 0{,}216 = 0{,}6^3
\;\Rightarrow\; \left(\frac{a}{L}\right)^{2/3} = 0{,}6^2 = 0{,}36 \\
y_{\mathrm m} &= 12{,}5\,(1 - 0{,}36)^{3/2} = 12{,}5 \cdot 0{,}64^{3/2}
= 12{,}5 \cdot 0{,}8^3 = 12{,}5 \cdot 0{,}512 = 6{,}4\,\text{m} \quad \checkmark \\
x_{\mathrm m} &= \sqrt[3]{2{,}7 \cdot 12{,}5^2} = \sqrt[3]{2{,}7 \cdot 156{,}25}
= \sqrt[3]{421{,}875} = 7{,}5\,\text{m} \quad \checkmark
\end{align}

Slajd uz to prikazuje graf $y(x)$ za te brojke: krivulja kreće od nule pri
$x = a = 2{,}7$, raste do vrha $y \approx 6{,}4$ pri $x = 7{,}5$, i pada natrag
na nulu pri $x = L = 12{,}5$ (tada su ljestve okomite uz zid).

\subsection{Optimalna ograda}
\str{52}

\textbf{Problem.} Površinu od 1 ha treba ograditi na način prema skici
(pravokutnik $x \times y$ s jednim unutarnjim pregradnim zidom duljine $y$).
Troškovi postavljanja ograde proporcionalni su duljini ograde. Odrediti
dimenzije $x$ i $y$ tako da troškovi budu minimalni.

\textbf{Ciljna funkcija je duljina ograde:}
\begin{equation}
L = 2x + 3y
\end{equation}

\emph{Odakle:} vanjski opseg je $2x + 2y$, a unutarnja pregrada je još jedan
komad duljine $y$; ukupno $2x + 3y$.

\textbf{Ograničenje je površina:}
\begin{equation}
A = x\,y = 10\,000\,\text{m}^2 \quad \to \quad h = xy - 10\,000 = 0
\end{equation}

(1 ha $=$ 10\,000 m\textsuperscript{2}. Zapis $h = xy - 10\,000 = 0$ je upravo
standardna forma ograničenja jednakosti iz
odjeljka~\ref{sec:standardna-forma}.)

\textbf{Uvrštenjem (2) u (1)}, tj.\ $y = 10\,000/x$, \textbf{ciljna funkcija
postaje:}
\[
L = 2x + \frac{3 \cdot 10\,000}{x}
\quad \to \quad
\frac{\dd L}{\dd x} = 2 - \frac{30\,000}{x^2} = 0
\]

\textbf{pa je:}
\[
x_{\text{opt}} = 100\sqrt{\frac{3}{2}}, \qquad
y_{\text{opt}} = 100\sqrt{\frac{2}{3}}, \qquad
\left(\frac{x}{y}\right)_{\text{opt}} = \frac{3}{2}
\]

\emph{Raspis:}
\begin{align}
2 &= \frac{30\,000}{x^2} \;\Rightarrow\; x^2 = 15\,000
\;\Rightarrow\; x = \sqrt{15\,000} = \sqrt{10\,000 \cdot 1{,}5}
= 100\sqrt{\tfrac{3}{2}} \approx 122{,}47\,\text{m} \\
y &= \frac{10\,000}{x} = \frac{10\,000}{100\sqrt{3/2}} = \frac{100}{\sqrt{3/2}}
= 100\sqrt{\tfrac{2}{3}} \approx 81{,}65\,\text{m} \\
\frac{x}{y} &= \frac{100\sqrt{3/2}}{100\sqrt{2/3}} = \sqrt{\frac{3/2}{2/3}}
= \sqrt{\frac{9}{4}} = \frac{3}{2} \quad \checkmark
\end{align}

\textbf{Interpretacija:} omjer stranica nije $1{:}1$ (kao kod običnog
pravokutnika bez pregrade) nego $3{:}2$ --- jer pregradni zid „poskupljuje”
smjer $y$, pa se isplati taj smjer skratiti.

% =====================================================================
\section{Vrste optimizacije konstrukcija}
\str{53--56}

\subsection{Dimenzionalna optimizacija (\emph{size optimization})}
\str{53}

Optimiranje \textbf{dimenzija} objekata \textbf{poznate strukture i oblika}.

Primjeri:
\begin{itemize}
  \item dimenzije elemenata nosivih konstrukcija (presjeci štapova, debljine
        limova, \dots)
  \item dimenzije članova mehanizama
\end{itemize}

\subsection{Optimiranje oblika (\emph{shape optimization})}
\str{54}

Osnovni \textbf{oblik} objekta mijenja se postupkom optimizacije ovisno o
zadanim parametrima problema (alat je često \textbf{FEM} analiza --- metoda
konačnih elemenata).

\subsection{Topološka optimizacija (\emph{topology optimization})}
\str{55--56}

Traži se optimalna \textbf{topologija (struktura)} konstrukcije u zadanoj
domeni.

\begin{napomena}
Na slajdovima 55 i 56 topološka optimizacija je numerirana kao „2.” (očit
tipfeler u prezentaciji --- dvije uzastopne stavke nose broj 2). Sadržajno je
to \textbf{treća} vrsta.
\end{napomena}

\textbf{Razlika u tri riječi:} \emph{dimenzionalna} mijenja \textbf{brojeve}
(koliko je štap debeo), \emph{oblikovna} mijenja \textbf{konture} (kako je rub
zaobljen), \emph{topološka} mijenja \textbf{raspored materijala} (gdje uopće ima
materijala, a gdje su rupe). Topološka je najslobodnija i najteža.

% =====================================================================
\section{Postavka problema}
\str{57--58}

\begin{teorem}[Ključna poruka slajda 57]
Pravilna definicija, formulacija i točna matematička postavka optimizacijskog
problema je često \textbf{najveći korak} prilikom rješavanja problema.

Iz formulacije problema treba raspoznati o \textbf{kojoj se vrsti optimizacije
radi}.
\end{teorem}

\subsection{Jednostavan primjer koji to ilustrira}
\str{57--58}

\textbf{Problem.} Površinu uz postojeću prirodnu ogradu (zid, stijena i sl.)
treba ograditi s raspoloživom duljinom žičane ograde $L$ tako da ograđena
površina bude \textbf{maksimalna}.

Na slajdu su prikazana \textbf{tri oblika}:
\begin{itemize}
  \item \textbf{a)} pravokutnik naslonjen na zid (jedna stranica $x$ uz zid,
        dvije stranice $y$),
  \item \textbf{b)} polukrug polumjera $r$ naslonjen na zid,
  \item \textbf{c)} pola pravilnog mnogokuta ($n$-terokut, konkretno $n = 8$)
        naslonjenog na zid.
\end{itemize}

\textbf{Poanta (sa slajda):}

\begin{quote}
Iz gornje formulacije nije jasno da li se oblik ograde može odabrati po volji
ili se postavljaju i neki drugi zahtjevi. Svaki od pokazanih oblika daje za istu
duljinu ograde svoje rješenje, pri čemu se oblikom b) ograđuje \textbf{najveća}
površina, a oblikom a) \textbf{najmanja}. Međutim, ako se radi o ogradi pašnjaka
koju treba seliti duž zida, tada je oblik a) u „tehnološkoj” prednosti, jer se
iskoristi sva raspoloživa površina.
\end{quote}

\textbf{Zaključak:} ista rečenica („ogradi maksimalnu površinu”) daje
\textbf{tri različita optimizacijska problema}, ovisno o tome što smo prešutno
pretpostavili. Zato je formulacija najvažniji korak.

\subsection{Brojčana usporedba triju oblika, $L = 500$ m}
\str{58}

\subsubsection*{Oblik b) --- polukrug}
\[
L = r\pi \;\to\; A = \frac{\pi r^2}{2} = \frac{L^2}{2\pi} = \frac{500^2}{2\pi}
= 39\,788{,}7\,\text{m}^2
\]

\emph{Raspis:} iz $L = r\pi$ slijedi $r = L/\pi$, pa je
$A = \tfrac{\pi}{2}\cdot\tfrac{L^2}{\pi^2} = \tfrac{L^2}{2\pi}
= \tfrac{250\,000}{6{,}2832} = 39\,788{,}7$ m\textsuperscript{2} \checkmark

\subsubsection*{Oblik c) --- polovica pravilnog $n$-terokuta, $n = 8$ (osmerokut)}
\[
L = \frac{n}{2}a \;\to\; \varphi = \frac{2\pi}{n}; \qquad
a = 2r\sin\frac{\varphi}{2} = 2r\sin\frac{\pi}{n}
\]

($a$ je duljina jedne stranice, $\varphi$ središnji kut nad jednom stranicom,
$r$ polumjer opisane kružnice. Uz zid je polovica mnogokuta, pa ograda ima $n/2$
stranica.)

\[
r = \frac{L}{n \cdot \sin\frac{\pi}{n}}
\]

\emph{Raspis:} iz $L = \tfrac{n}{2}a = \tfrac{n}{2}\cdot 2r\sin\tfrac{\pi}{n}
= n\,r\sin\tfrac{\pi}{n}$ slijedi $r = \tfrac{L}{n\sin(\pi/n)}$ \checkmark

\[
A = \frac{1}{2}a\,r\cos\frac{\varphi}{2}\cdot\frac{n}{2}
= \frac{n}{4}\cdot 2\frac{L}{n}\cdot\frac{L}{n\cdot\sin\frac{\pi}{n}}
  \cos\frac{\pi}{n}
= \frac{L^2}{2\,n \cdot \tan\frac{\pi}{n}}
\]

\[
A = \frac{500^2}{2 \cdot 8 \cdot \tan\frac{\pi}{8}} = 37\,722{,}1\,\text{m}^2
\qquad \left(\text{za } n \to \infty, \;
\lim_{n\to\infty}\left(n\cdot\tan\frac{\pi}{n}\right) = \pi\right)
\]

\emph{Provjera brojki:} $\tan(\pi/8) = \tan 22{,}5^\circ = 0{,}41421$, pa je
$2 \cdot 8 \cdot 0{,}41421 = 6{,}6274$ i
$A = 250\,000/6{,}6274 = 37\,722$ m\textsuperscript{2} \checkmark

\emph{Značenje granice:} kad $n \to \infty$, $n\tan(\pi/n) \to \pi$, pa formula
za mnogokut prelazi u $A \to L^2/(2\pi)$ --- točno rezultat polukruga b).
Mnogokut s mnogo stranica je praktički polukrug, što potvrđuje da je b)
najbolji oblik.

\subsubsection*{Oblik a) --- pravokutnik}
\[
L = x + 2y \;\to\; x = L - 2y \quad (\text{ograničenje jednakosti})
\]
\[
A = xy = (L - 2y)\,y; \qquad
\frac{\dd A}{\dd y} = L - 4y = 0 \;\to\;
y = \frac{L}{4} = 125\,\text{m}; \quad x = \frac{L}{2} = 250\,\text{m}
\]

\emph{Raspis derivacije:} $A = Ly - 2y^2$, pa je
$\tfrac{\dd A}{\dd y} = L - 4y$ \checkmark

\[
A = x\,y = 250 \cdot 125 = 31\,250\,\text{m}^2 \qquad (x : y = 2 : 1)
\]

\textbf{Rangiranje:} b) $39\,788{,}7$ $>$ c) $37\,722{,}1$ $>$ a) $31\,250$
m\textsuperscript{2}. Razlika između najboljeg i najlošijeg oblika je oko
\textbf{27\,\%} --- a sve uz istu duljinu ograde. Odabir \emph{strukture}
(oblika) donosi veći dobitak nego optimiranje \emph{dimenzija} unutar zadane
strukture. To je uvod u sljedeći odjeljak.

% =====================================================================
\section{Projektiranje kao proces}
\str{59--77}

\subsection{Slijed projektiranja}
\str{59}

Projektiranje tehničkog sustava ostvaruje se u \textbf{dva koraka}:

\begin{enumerate}
  \item \textbf{izbor strukture} sustava, odnosno \textbf{strukturalne sinteze};
  \item \textbf{izbor brojčanih vrijednosti parametara} odabrane strukture
        (sinteza parametara ili \textbf{dimenzionalna sinteza}).
\end{enumerate}

\begin{quote}
\textbf{Strukturalna sinteza} jest (glavni) zadatak \textbf{stvaralačkog
djelovanja inženjera} koji je općenito teško ili nemoguće automatizirati
(algoritmizirati).
\end{quote}

To se izravno poklapa s primjerom ograde: odabir između oblika a), b), c) je
strukturalna sinteza (kreativni čin), a računanje $x = 250$, $y = 125$ je
dimenzionalna sinteza (algoritamski dio) --- i baš je taj drugi dio ono što ovaj
kolegij formalizira.

\subsection{Strukturalna sinteza}
\str{60}

Struktura projektiranog tehničkog sustava ili objekta (\textbf{topologija
konstrukcije}) određena je poznatim \textbf{skupom elemenata strukture} i
\textbf{vezama među njima}.

\textbf{Poteškoće koje se pojavljuju u formalizaciji problema strukturalne
sinteze:}
\begin{itemize}
  \item nemjerljive karakteristike elemenata strukture;
  \item formalno nedefinirane veze između pojedinih elemenata i zahtjeva;
  \item kvalitativni kriteriji;
  \item neformalan opis funkcioniranja projektiranog objekta;
  \item zahtjevi koji se odnose na funkcioniranje objekta.
\end{itemize}

\textbf{Temelj sinteze tehničkih sustava jest:} skup znanja i iskustva
projektanata, intuicija, kreativnost, literatura, izvedeni projekti, sheme
mogućih rješenja, metode projektiranja, računalna i programska podrška.

Opis elemenata strukture, opis veza među njima kao i formiranje matematičkog
modela strukture definira se pomoću prikladnog (konačnog) skupa veličina: to su
\textbf{projektni parametri}.

\subsection{Cilj projektiranja?}
\str{61--67}

\begin{enumerate}
  \item Gradnja sustava s besprijekornim funkcioniranjem se \textbf{danas
        podrazumijeva}.
  \item U uvjetima konkurencije od presudne je važnosti projektirati
        \textbf{najbolji ili što bolji} sustav (objekt ili proizvod).
  \item Kriteriji po kojima je neki sustav najbolji su različiti, primjerice:
        \textbf{cijena, težina, efikasnost, oblik, praktičnost, kompleksnost,
        originalnost, pouzdanost, trajnost} i sl.
  \item \textbf{Izbor dominantnog kriterija} predstavlja \textbf{stratešku
        odluku} na početku samog procesa projektiranja. Kvaliteta rješenja
        ocjenjuje se \textbf{kvantitativnom} usporedbom odabranih kriterija,
        koji stoga moraju biti \textbf{funkcijom projektnih parametara}.
  \item Izbor projektnih parametara najčešće \textbf{nije slobodan}, već njihove
        veličine moraju zadovoljiti neke unaprijed zadane uvjete ili
        \textbf{ograničenja}. Također trebaju biti ispunjeni uvjeti
        \textbf{međusobne ovisnosti} parametara koji slijede iz strukture
        projektiranog sustava.
  \item Projektiranje je stoga \textbf{iterativni proces} određivanja strukture
        i projektnih parametara projektiranog sustava, konstrukcije ili
        proizvoda.
  \item Konačnim projektnim rješenjem trebaju se ispuniti \textbf{svi} kriteriji
        i tehnički zahtjevi projekta.
\end{enumerate}

Točka 4 je zapravo definicija \textbf{funkcije cilja} rečena inženjerskim
jezikom („kriterij mora biti funkcija projektnih parametara”), a točka 5
definicija \textbf{ograničenja}. Time se zatvara krug s
odjeljkom~\ref{sec:standardna-forma}.

\subsection{Proces projektiranja}
\str{68}

Prema primijenjenim metodama i projektnim zahtjevima proces projektiranja može
biti:

\textbf{a) Proces konvencionalnog projektiranja.} Osnovni cilj: besprijekorno
funkcioniranje sustava ili proizvoda u okviru uvjeta postavljenih projektnim
zadatkom (\emph{funkcionalno projektiranje}). Ovaj proces baziran je na
\textbf{iskustvu i intuiciji}.

\textbf{b) Proces optimalnog projektiranja.} Pored besprijekornog
funkcioniranja sustava ili proizvoda, \textbf{jedan ili više projektnih
kriterija treba biti ispunjen na optimalan način}. Postupak optimalnog
projektiranja zahtijeva \textbf{strogu matematičku formulaciju problema} ---
postavku matematičkog modela ili optimizacijskog problema.

\subsection{Elementi optimizacijskog modela: projektni parametri i varijable}
\str{69--70}

\begin{definicija}[Projektni parametri]
\textbf{Projektni parametri} su veličine pomoću kojih se \textbf{jednoznačno
definiraju dimenzije i svojstva} konstrukcije ili projektiranog sistema.
\end{definicija}

Projektni parametri dijele se na \textbf{konstantne projektne parametre} i
\textbf{projektne varijable}.

\textbf{Konstantni projektni parametri} se u postupku optimizacije \textbf{ne
mijenjaju}, a mogu se podijeliti na dvije vrste:

\begin{itemize}
  \item \textbf{taktički projektni parametri ($z$)} --- zadani su
        \textbf{projektnim zadatkom} i ne mogu se mijenjati u cijelom procesu
        projektiranja. Oni su konstante strukturalne i dimenzionalne sinteze
        konstrukcije, tj.\ predstavljaju \textbf{taktičke performanse}
        projektiranog sustava.
  \item \textbf{tehnički projektni parametri ($y$)} --- projektne su konstante
        \textbf{po zamisli projektanta}, i to prvenstveno iz
        \textbf{konstrukcijskih ili tehnoloških} razloga.
\end{itemize}

\textbf{Hijerarhija sa slajda 70:}

\begin{center}
\begin{tabular}{c}
\toprule
\textbf{PROJEKTNI PARAMETRI} \\
\midrule
\begin{tabular}{p{5.0cm}|p{7.0cm}}
\centering \textbf{PROJEKTNE VARIJABLE} & \centering\textbf{PROJEKTNE KONSTANTE} \tabularnewline
\centering $x$ & \centering $y,\; z$ \tabularnewline
 & \begin{tabular}{p{3.0cm}p{3.0cm}}
   \textbf{TEHNIČKI} $y$ & \textbf{TAKTIČKI} $z$ \\
   (odabrani od projektanta) & (zadani projektnim zadatkom)
   \end{tabular}
\end{tabular} \\
\bottomrule
\end{tabular}
\end{center}

\textbf{Dodatne tvrdnje sa slajda 70:}
\begin{itemize}
  \item Tehnički parametri su konstante \textbf{samo kod dimenzionalne sinteze
        već poznate strukture (topologije)} objekta.
  \item Projektant može mijenjati broj tehničkih parametara njihovim
        \textbf{prebacivanjem u projektne varijable i obrnuto}.
  \item Projektni parametri koji se \textbf{mijenjaju} u postupku sinteze
        nazivaju se \textbf{projektnim varijablama ($x$)}.
  \item \textbf{Suma broja projektnih varijabli i tehničkih parametara} u
        postupku optimizacije iste strukture sistema je \textbf{konstantna}.
  \item Ispravan izbor projektnih parametara i odgovarajuća matematička
        formulacija problema glavne su poteškoće i preduvjeti uspješnog postupka
        optimalnog projektiranja.
\end{itemize}

\begin{napomena}[veza sa standardnom formom]
Vektor $x$ iz standardne forme $=$ projektne varijable. Parametri $y$ i $z$ ne
pojavljuju se kao nepoznanice, nego kao \textbf{brojevi} u izrazima za $f$, $h$
i $g$. U primjeru ograde: $x$ i $y$ (dimenzije) su projektne varijable, a
$10\,000$ m\textsuperscript{2} je taktički parametar (zadan zadatkom).
\end{napomena}

\subsection{Optimizacijski model --- pravila izbora parametara}
\str{71--77}

Kod izbora projektnih parametara važno je sljedeće:

\begin{enumerate}
  \item Projektant treba dobro \textbf{proučiti problem} i matematički ga
        opisati pomoću projektnih varijabli;
  \item Projektni parametri trebaju biti \textbf{nezavisni};
  \item Problem treba opisati sa \textbf{što manje} projektnih parametara;
  \item Matematičku formulaciju problema treba u \textbf{početnoj fazi} provesti
        sa \textbf{što manjim brojem tehničkih parametara}; to je važno radi
        fleksibilnosti algoritma i otklanjanja eventualnih pogrešnih procjena o
        utjecaju nekih tehničkih parametara na rezultat. U \textbf{kasnijoj
        fazi} mogu se neke od varijabli odabrati kao konstante, i to:
        \begin{itemize}
          \item iz \textbf{konstrukcijskih ili tehnoloških} razloga;
          \item jer je utvrđeno da \textbf{ne utječu bitno} na konačni rezultat
                optimizacije;
          \item jer je utvrđeno da optimalno rješenje odgovara \textbf{uvijek
                vrijednosti neke varijable na rubu ograničenja}.
        \end{itemize}
\end{enumerate}

\begin{napomena}
\textbf{Zašto pravilo 2 (nezavisnost)?} Ako su dvije varijable povezane (npr.\
uvedemo i polumjer $r$ i promjer $d = 2r$), ograničenje $d = 2r$ postaje suvišna
jednadžba, a problem numerički degenerira (rješenje nije jedinstveno u
parametrima). Bolje je odabrati samo jednu od njih.

\textbf{Zašto pravilo 3 (što manje parametara)?} Broj varijabli $n$ izravno
određuje dimenziju prostora u kojem tražimo optimum i time složenost
rješavanja.

Zadnja alineja pravila 4 najavljuje pojam \textbf{aktivnog ograničenja} iz
poglavlja~\ref{ch:uvjeti}: ako znamo da optimum uvijek leži na rubu, možemo tu
nejednakost zamijeniti jednakošću i eliminirati varijablu.
\end{napomena}

% =====================================================================
\section{Neki korisni „trikovi”}
\label{sec:trikovi}
\str{79}

Ovaj slajd zatvara predavanje 1 i ponavlja se na početku predavanja 2.

\subsection{Minimizacija $\Leftrightarrow$ maksimizacija}

\begin{equation}
\begin{array}{ccc}
\begin{array}{l}
\max_x f(x) \\ \text{uz ogr.} \quad x \in \Feas
\end{array}
& \Longleftrightarrow &
\begin{array}{l}
\min_x -f(x) \\ \text{uz ogr.} \quad x \in \Feas
\end{array}
\end{array}
\end{equation}

\textbf{Zašto to vrijedi.} Ako $f$ postiže svoju \textbf{najveću} vrijednost u
točki $x\opt$, onda $-f$ u istoj toj točki postiže svoju \textbf{najmanju}
vrijednost (jer množenje s $-1$ „okreće” os vrijednosti). Dakle:

\begin{itemize}
  \item \textbf{argument} je isti: $\argmax f = \argmin (-f)$,
  \item \textbf{vrijednosti} su suprotne: $\max f = -\min(-f)$ --- na to treba
        paziti kad izvještavate optimalnu vrijednost!
\end{itemize}

Zbog ovog trika \textbf{cijela teorija se može razviti samo za minimizaciju},
što i činimo u standardnoj formi iz odjeljka~\ref{sec:standardna-forma}.

\subsection{Provjera je li dozvoljeni skup prazan (problem dopustivosti)}

\textbf{Pitanje:} Je li skup $\Feas = \{\, x \mid g_i(x) \le 0,\;
i = 1,\dots,p \,\}$ prazan? Ako nije, kako naći neki $x \in \Feas$?

\textbf{Rješenje:} riješi optimizacijski problem

\begin{equation}
\label{eq:dopustivost}
\begin{aligned}
\min_{x,\,t} \quad & t \\
\uzogr & g_i(x) \le t, \quad i = 1,\dots,p
\end{aligned}
\end{equation}

Neka je $t\opt, x\opt$ rješenje. \textbf{Ako je $t\opt \le 0$, tada $\Feas$ nije
prazan skup i $x\opt \in \Feas$.}

\textbf{Detaljno objašnjenje} (jer je ideja lukavija nego što izgleda):

\begin{itemize}
  \item Uvodimo \textbf{novu, pomoćnu varijablu} $t \in \R$. Sada minimiziramo
        po paru $(x, t)$, dakle problem ima $n + 1$ varijablu umjesto $n$.
  \item Za \emph{bilo koji} $x$ uvijek možemo naći dovoljno velik $t$ da svi
        uvjeti $g_i(x) \le t$ budu zadovoljeni (npr.\ $t = \max_i g_i(x)$). Zato
        \textbf{novi problem nikad nije nedopustiv} --- uvijek ima rješenje, za
        razliku od polaznog.
  \item Funkcija cilja je samo $t$: guramo $t$ što niže. Za fiksni $x$ najmanji
        dopušteni $t$ je točno $\max_i g_i(x)$ --- „najgore prekršeno
        ograničenje”. Dakle problem \eqref{eq:dopustivost} zapravo rješava
        \[
        \min_x \; \max_{i} \; g_i(x)
        \]
        tj.\ traži $x$ koji \textbf{najmanje krši najgore ograničenje}.
  \item \textbf{Ako je $t\opt \le 0$:} tada je $g_i(x\opt) \le t\opt \le 0$ za
        sve $i$, dakle $x\opt$ zadovoljava \textbf{sva} izvorna ograničenja
        $\to$ $x\opt \in \Feas$ i skup nije prazan. Bonus: dobili smo i
        konkretnu dopuštenu točku.
  \item \textbf{Ako je $t\opt > 0$:} ne postoji nijedan $x$ za koji su sva
        $g_i(x) \le 0$ (jer bi za takav $x$ vrijednost $t = 0$ bila dopuštena i
        dala manji cilj od $t\opt$) $\to$ $\Feas$ \textbf{je prazan}.
\end{itemize}

\begin{napomena}[faza I]
Ovaj se postupak zove \textbf{faza I} (engl.\ \emph{phase I}) i standardni je
prvi korak svakog ozbiljnog optimizacijskog algoritma: prije nego što tražimo
\emph{najbolju} dopuštenu točku, moramo naći \emph{bilo koju}. Uočite i
tehniku: \textbf{uvođenje pomoćne varijable koja „gura” min--max u obični min}
--- istu ćemo tehniku sresti kod norma (Vježbe 2, poglavlje~\ref{ch:definicije})
i kod robusnog optimiranja (poglavlje~\ref{ch:robusno}).
\end{napomena}

% =====================================================================
\section{Sažetak poglavlja}

\textbf{Što morate znati napamet nakon ovog poglavlja:}

\begin{enumerate}
  \item \textbf{Standardna forma:} $\min f(x)$ uz $h(x) = 0$, $g(x) \le 0$; što
        je $x$, $f$, $h$, $g$, i što je dozvoljeni skup
        $\Feas = \{x \mid h(x)=0,\, g(x)\le 0\}$.
  \item \textbf{Definicija optimalnog rješenja} $x\opt$ (dopustivo \emph{i}
        najmanji $f$).
  \item \textbf{Nivo krivulja} $\{x \mid f(x) = c\}$ i kako se iz slike nivo
        krivulja $+$ dozvoljenog skupa očita optimum
        (odjeljak~\ref{sec:uvodni-primjer}).
  \item \textbf{Prevođenje u standardnu formu:} premještanje konstante na lijevu
        stranu, okretanje „$\ge$” množenjem s $-1$,
        $\max f \Leftrightarrow \min(-f)$.
  \item \textbf{Recept klasične analize} (eliminacija ograničenja jednakosti
        uvrštavanjem $+$ $F' = 0$ $+$ provjera s $F''$) --- i svijest da on ne
        funkcionira općenito, zbog čega nam trebaju
        poglavlja~\ref{ch:uvjeti}~i~\ref{ch:dualnost}.
  \item \textbf{Problem dopustivosti} \eqref{eq:dopustivost} i njegova
        interpretacija.
\end{enumerate}

\textbf{Terminologija (hrvatski $\leftrightarrow$ engleski):}

\begin{center}
\begin{tabular}{ll}
\toprule
\textbf{Hrvatski} & \textbf{Engleski} \\
\midrule
projektne varijable / varijable odluke & decision variables \\
funkcija cilja / ciljna funkcija & objective function \\
ograničenja jednakosti & equality constraints \\
ograničenja nejednakosti & inequality constraints \\
dozvoljeni skup / dozvoljeno područje & feasible set \\
nivo krivulja & level curve / contour \\
dimenzionalna optimizacija & size optimization \\
optimiranje oblika & shape optimization \\
topološka optimizacija & topology optimization \\
\bottomrule
\end{tabular}
\end{center}
